\section{Sơ đồ Thực thể Kết hợp (ERD)}

Dựa trên danh sách các thực thể và mối quan hệ đã phân tích, sơ đồ ERD ở mức quan niệm được thiết kế theo ký pháp chuẩn Chen (Thực thể dạng hình chữ nhật, Mối kết hợp dạng hình thoi, Thuộc tính dạng hình bầu dục). Sơ đồ này độc lập với hệ quản trị CSDL vật lý.

% =========================================================
% TÙY CHỌN 1: SƠ ĐỒ TIKZ FULL THUỘC TÍNH (BẢN XEM TRƯỚC ĐÃ TỐI ƯU TỶ LỆ)
% =========================================================
\begin{figure}[H]
    \centering
    \resizebox{\textwidth}{!}{
    \begin{tikzpicture}[>=stealth, thick]
        
        % 1. ĐỊNH NGHĨA LẠI STYLE: Hình chữ nhật (nhỏ gọn), Hình thoi (to và cân đối)
        \tikzstyle{entity} = [rectangle, draw=black, fill=blue!10, minimum width=2.2cm, minimum height=0.7cm, text centered, font=\footnotesize\bfseries]
        \tikzstyle{weak_entity} = [rectangle, draw=black, double, double distance=1.5pt, fill=blue!10, minimum width=2.4cm, minimum height=0.7cm, text centered, font=\footnotesize\bfseries]
        \tikzstyle{relationship} = [diamond, draw=black, fill=orange!20, minimum width=2.8cm, minimum height=1.6cm, text centered, font=\footnotesize\itshape]
        \tikzstyle{weak_relationship} = [diamond, draw=black, double, double distance=1.5pt, fill=orange!20, minimum width=2.8cm, minimum height=1.6cm, text centered, font=\footnotesize\itshape]
        \tikzstyle{attribute} = [ellipse, draw=black, fill=white, inner sep=2pt, minimum width=1.5cm, minimum height=0.6cm, text centered, font=\scriptsize]

        % 2. CÁC THỰC THỂ (Kéo giãn khoảng cách tọa độ để không bị rối)
        \node[entity] (tenant) at (0, 12) {Tenant};
        \node[entity] (merchant) at (0, 8) {Merchant};
        \node[entity] (user) at (-8, 8) {User};
        \node[entity] (role) at (-12, 11) {Role}; % Thêm thực thể ROLE
        \node[entity] (campaign) at (8, 8) {Campaign};
        \node[entity] (rule) at (15, 8) {RewardRule};
        \node[entity] (customer) at (-8, 3) {Customer};
        \node[entity] (credit) at (0, 3) {Credit};
        \node[entity] (api) at (8, 4.5) {APICredential};
        \node[entity] (transaction) at (0, -3) {CreditTransaction};

        % 3. MỐI KẾT HỢP
        \node[relationship] (capphat) at (0, 10) {Cấp phát};
        \node[relationship] (quanly) at (-4, 8) {Quản lý};
        \node[relationship] (phanquyen) at (-10, 9.5) {Phân quyền}; % Quan hệ Role - User
        \node[relationship] (sếpnhanvien) at (-10.5, 6.5) {Sếp-NV}; % Quan hệ tự tham chiếu User - User
        \node[relationship] (khoitao) at (4, 8) {Khởi tạo};
        \node[relationship] (baogom) at (11.5, 8) {Bao gồm};
        \node[relationship] (phathanh) at (0, 5.5) {Phát hành};
        \node[relationship] (sohuu) at (-4, 3) {Sở hữu};
        \node[relationship] (taokhoa) at (4, 6.25) {Tạo khóa};
        \node[relationship] (ghinhan) at (0, 0) {Ghi nhận};
        \node[relationship] (apdung) at (8, 0) {Áp dụng}; 

        % 4. NỐI MỐI KẾT HỢP (Nhãn 1-N)
        \draw (tenant) -- node[right] {1} (capphat) -- node[right] {N} (merchant);
        \draw (user) -- node[above] {N} (quanly) -- node[above] {1} (merchant);
        \draw (role) -- node[left] {1} (phanquyen) -- node[right] {N} (user);
        \draw (user.west) -- node[above] {0..1} (sếpnhanvien);
        \draw (sếpnhanvien.south) |- node[below, pos=0.75] {N} (user.south west); % Vòng tự tham chiếu ManagerID
        \draw (merchant) -- node[above] {1} (khoitao) -- node[above] {N} (campaign);
        \draw (campaign) -- node[above] {1} (baogom) -- node[above] {N} (rule);
        \draw (merchant) -- node[right] {1} (phathanh) -- node[right] {N} (credit);
        \draw (customer) -- node[above] {1} (sohuu) -- node[above] {N} (credit);
        \draw (merchant) -- node[above, pos=0.8] {1} (taokhoa) -- node[above, pos=0.8] {N} (api);
        \draw (credit) -- node[right] {1} (ghinhan) -- node[right] {N} (transaction);
        \draw (campaign) -- node[right, pos=0.4] {0..1} (apdung);
        \draw (apdung.south) |- node[above, pos=0.9] {N} (transaction.east);

        % 5. VẼ THUỘC TÍNH (Tỏa đều ra không đè nét)
        % Tenant
        \node[attribute] (t1) at (-2.5, 13) {\underline{TenantID}};
        \node[attribute] (t2) at (0, 13.5) {TenantName};
        \node[attribute] (t3) at (2.5, 13) {SystemConfig};
        \draw (tenant) -- (t1); \draw (tenant) -- (t2); \draw (tenant) -- (t3);

        % Merchant
        \node[attribute] (m1) at (-2.5, 9.5) {\underline{MerchantID}};
        \node[attribute] (m2) at (2.5, 9.5) {BusinessName};
        \node[attribute] (m3) at (-2, 6.5) {Status};
        \draw (merchant) -- (m1); \draw (merchant) -- (m2); \draw (merchant) -- (m3);

        % Role
        \node[attribute] (role1) at (-14, 11.5) {\underline{RoleID}};
        \node[attribute] (role2) at (-14, 10.5) {RoleName};
        \draw (role) -- (role1); \draw (role) -- (role2);

        % User (cập nhật thêm RoleID và ManagerID)
        \node[attribute] (u1) at (-12.5, 8.5) {\underline{UserID}};
        \node[attribute] (u2) at (-10.5, 8.5) {Email};
        \node[attribute] (u3) at (-6.5, 9.5) {RoleID};
        \node[attribute] (u4) at (-10.5, 5.5) {PasswordHash};
        \node[attribute] (u5) at (-6.5, 6.5) {ManagerID};
        \draw (user) -- (u1); \draw (user) -- (u2); \draw (user) -- (u3); \draw (user) -- (u4); \draw (user) -- (u5);

        % Campaign
        \node[attribute] (c1) at (6, 9.5) {\underline{CampaignID}};
        \node[attribute] (c2) at (8, 9.5) {CampaignName};
        \node[attribute] (c3) at (10, 9.5) {Status};
        \node[attribute] (c4) at (10, 6.5) {StartDate};
        \node[attribute] (c5) at (6, 6.5) {EndDate};
        \draw (campaign) -- (c1); \draw (campaign) -- (c2); \draw (campaign) -- (c3); \draw (campaign) -- (c4); \draw (campaign) -- (c5);

        % RewardRule
        \node[attribute] (r1) at (14, 9.5) {\underline{RuleID}};
        \node[attribute] (r2) at (16, 9.5) {ActionType};
        \node[attribute] (r3) at (16.5, 6.5) {PointsMultiplier};
        \draw (rule) -- (r1); \draw (rule) -- (r2); \draw (rule) -- (r3);

        % APICredential
        \node[attribute] (a1) at (11, 5.5) {\underline{KeyID}};
        \node[attribute] (a2) at (11.5, 4.5) {ApiKey};
        \node[attribute] (a3) at (11, 3.5) {SecretHash};
        \draw (api) -- (a1); \draw (api) -- (a2); \draw (api) -- (a3);

        % Customer (Cập nhật chuẩn Walletless Auth)
        \node[attribute] (cus1) at (-10.5, 4.5) {\underline{CustomerID}};
        \node[attribute] (cus2) at (-8, 4.5) {Email};
        \node[attribute] (cus3) at (-5.5, 4.5) {WalletAddress};
        \draw (customer) -- (cus1); \draw (customer) -- (cus2); \draw (customer) -- (cus3);

        % Credit
        \node[attribute] (cre1) at (-2.5, 1.5) {\underline{CreditID}};
        \node[attribute] (cre2) at (2.5, 1.5) {AvailableBalance};
        \draw (credit) -- (cre1); \draw (credit) -- (cre2);

        % CreditTransaction
        \node[attribute] (tr1) at (-4, -4.5) {\underline{TransactionID}};
        \node[attribute] (tr2) at (-1.5, -5) {Amount};
        \node[attribute] (tr3) at (1.5, -5) {TransactionType};
        \node[attribute] (tr4) at (4, -4.5) {SuiTxDigest};
        \node[attribute] (tr5) at (-3.5, -1.5) {CreatedAt};
        \draw (transaction) -- (tr1); \draw (transaction) -- (tr2); \draw (transaction) -- (tr3); \draw (transaction) -- (tr4); \draw (transaction) -- (tr5);

    \end{tikzpicture}
    }
    \caption{Sơ đồ Thực thể Kết hợp (ERD) mức quan niệm hệ thống OctaPoint CaaS}
    \label{fig:erd_tikz}
\end{figure}
% =========================================================
% TÙY CHỌN 2: CHÈN ẢNH DRAW.IO (KHI NÀO VẼ XONG THÌ DÙNG)
% Hướng dẫn:
% 1. Xóa (hoặc bôi đen nhấn Ctrl + /) toàn bộ đoạn code \begin{figure} của Tùy chọn 1 ở trên.
% 2. Xóa các dấu % ở đầu các dòng code của phần Tùy chọn 2 bên dưới để kích hoạt.
% 3. Lưu ảnh ERD tải về từ draw.io vào thư mục images/ với tên erd_quan_niem.png
% =========================================================

% \begin{figure}[H]
%     \centering
%     \includegraphics[width=0.95\textwidth]{images/erd_quan_niem.png}
%     \caption{Sơ đồ Thực thể Kết hợp (ERD) mức quan niệm hệ thống OctaPoint CaaS}
%     \label{fig:erd_drawio}
% \end{figure}
CreditTransaction có khóa TransactionID riêng nên là thực thể mạnh, dù phụ thuộc tồn tại vào Credit. Quan hệ Ghi nhận không phải quan hệ định danh của thực thể yếu. Các nhãn 0..1 biểu diễn người quản lý và chiến dịch tùy chọn; mỗi cặp Customer--Merchant có tối đa một Credit.
